\section{Introduction}
The sequence-to-sequence framework has demonstrated success in
natural-language sequence
transduction tasks such as machine translation.
More recently,
neural techniques have been applied to do single-document,
abstractive (paraphrasing)
text summarization of news
articles (\cite{Rush2015}, \cite{nallapati2016abstractive}). In this prior work, the input to
supervised models ranged from the first sentence to the entire
text of an article, and they are trained end-to-end to predict reference summaries.
Doing this end-to-end requires a significant number of parallel
article-summary pairs since language understanding
is a pre-requisite to generate fluent summaries.

In contrast, we consider the task of multi-document summarization, where
the input is a collection of related documents from which a summary is
distilled. Prior work has focused on extractive summarization,
which select sentences or phrases from the input to form the summaries, rather
than generating new text.
There has been limited application of abstractive neural methods and one
possible reason is the paucity of large, labeled datasets.

In this work,
we consider English Wikipedia as a supervised machine learning task
for multi-document summarization where the input is comprised of a Wikipedia
topic (title of article) and a collection of non-Wikipedia reference documents,
and the target is the Wikipedia article text.
We describe the first attempt to abstractively
generate the first section, or \textit{lead}, of Wikipedia articles
conditioned on reference text. In addition to
running strong baseline models on the task, we modify the Transformer architecture 
\citep{vaswani2017attention} to only consist of a decoder,
which performs better in the case of
longer input sequences compared to recurrent neural network (RNN) and Transformer encoder-decoder models.
Finally we show our modeling improvements allow us to generate entire Wikipedia articles.
